\textbf{Phần 1. Trắc nghiệm lựa chọn}
\setcounterpageref{ex}{0}

\begin{ex} %[3D1N08]
Số gồm $7$ trăm, $4$ chục và $2$ đơn vị là
\choice
{$724$}
{$472$}
{\True $742$}
{$247$}
\loigiai{
Số đó là $742$.\\
\textbf{Đáp án của câu hỏi là C}
}
\end{ex}

\begin{ex} %[3D1N08]
Kết quả của phép tính $4 \times 7$ là
\choice
{$24$}
{\True $28$}
{$32$}
{$21$}
\loigiai{
$4 \times 7 = 28$.\\
\textbf{Đáp án của câu hỏi là B}
}
\end{ex}

\begin{ex} %[3D1H08]
Kết quả của phép tính $32 : 4$ là
\choice
{\True $8$}
{$7$}
{$9$}
{$6$}
\loigiai{
$32 : 4 = 8$.\\
\textbf{Đáp án của câu hỏi là A}
}
\end{ex}

\begin{ex} %[3D1H08]
Tìm $x$ biết: $x + 145 = 600$.
\choice
{$445$}
{$465$}
{$745$}
{\True $455$}
\loigiai{
$x = 600 - 145 = 455$.\\
\textbf{Đáp án của câu hỏi là D}
}
\end{ex}

\begin{ex} %[3D1H08]
Tìm $y$ biết: $y - 230 = 370$.
\choice
{$500$}
{$140$}
{\True $600$}
{$610$}
\loigiai{
$y = 370 + 230 = 600$.\\
\textbf{Đáp án của câu hỏi là C}
}
\end{ex}

\begin{ex} %[3D1H08]
Tính giá trị biểu thức: $3 \times 8 + 25$.
\choice
{\True $49$}
{$47$}
{$51$}
{$45$}
\loigiai{
$3 \times 8 + 25 = 24 + 25 = 49$.\\
\textbf{Đáp án của câu hỏi là A}
}
\end{ex}

\begin{ex} %[3D1V08]
Một cuộn dây điện dài $100\text{ m}$, người ta đã cắt đi 4 đoạn, mỗi đoạn dài $5\text{ m}$. Cuộn dây còn lại số mét là
\choice
{$75\text{ m}$}
{$85\text{ m}$}
{$90\text{ m}$}
{\True $80\text{ m}$}
\loigiai{
Đã cắt: $5 \times 4 = 20\text{ m}$.\\
Còn lại: $100 - 20 = 80\text{ m}$.\\
\textbf{Đáp án của câu hỏi là D}
}
\end{ex}

\begin{ex} %[3D1V08]
Có 27 học sinh chia đều thành 3 tổ. Hỏi 2 tổ có bao nhiêu học sinh?
\choice
{$16$ học sinh}
{\True $18$ học sinh}
{$20$ học sinh}
{$15$ học sinh}
\loigiai{
Mỗi tổ: $27 : 3 = 9$ bạn. 2 tổ: $9 \times 2 = 18$ bạn.\\
\textbf{Đáp án của câu hỏi là B}
}
\end{ex}

\begin{ex} %[3D1V08]
Sắp xếp các số $567, 765, 657, 576$ theo thứ tự từ lớn đến bé:
\choice
{$567 > 576 > 657 > 765$}
{$765 > 657 > 567 > 576$}
{$657 > 765 > 576 > 567$}
{\True $765 > 657 > 576 > 567$}
\loigiai{
Thứ tự giảm dần: $765 > 657 > 576 > 567$.\\
\textbf{Đáp án của câu hỏi là D}
}
\end{ex}

\begin{ex} %[3D1V08]
Tính chu vi hình tứ giác có các cạnh đều bằng $5\text{ cm}$.
\choice
{$15\text{ cm}$}
{$25\text{ cm}$}
{\True $20\text{ cm}$}
{$10\text{ cm}$}
\loigiai{
$5 \times 4 = 20\text{ cm}$.\\
\textbf{Đáp án của câu hỏi là C}
}
\end{ex}

\begin{ex} %[3D1V08]
Tìm $x$ biết: $x : 3 = 10 - 2$.
\choice
{\True $24$}
{$21$}
{$27$}
{$18$}
\loigiai{
$x : 3 = 8 \Rightarrow x = 8 \times 3 = 24$.\\
\textbf{Đáp án của câu hỏi là A}
}
\end{ex}

\begin{ex} %[3D1C08]
Một giá sách có 2 ngăn, ngăn trên có 45 cuốn sách. Nếu chuyển 5 cuốn từ ngăn trên xuống ngăn dưới thì hai ngăn bằng nhau. Hỏi lúc đầu ngăn dưới có bao nhiêu cuốn sách?
\choice
{$30$ cuốn}
{\True $35$ cuốn}
{$40$ cuốn}
{$25$ cuốn}
\loigiai{
Sau khi chuyển ngăn trên còn: $45 - 5 = 40$ cuốn $\Rightarrow$ ngăn dưới lúc sau có 40 cuốn.\\
Lúc đầu ngăn dưới có: $40 - 5 = 35$ cuốn sách.\\
\textbf{Đáp án của câu hỏi là B}
}
\end{ex}

\textbf{Phần 2. Trắc nghiệm đúng sai}
\setcounterpageref{ex}{0}

\begin{ex} %[3D1H08]
Xét tính đúng sai của các biểu thức toán học:
\choiceTF
{\True $300 + 400 + 50 = 750$.}
{\True $4 \times 8 = 32$.}
{$24 : 3 = 9$.}
{\True $24 : 3 = 8$.}
\loigiai{
\begin{itemize}
    \item $300 + 400 + 50 = 750$: Đúng.
    \item $4 \times 8 = 32$: Đúng.
    \item $24 : 3 = 8 \neq 9$. Ý 3 Sai, Ý 4 Đúng.
\end{itemize}
\begin{itemize}
    \item a) \textbf{Đúng}.
    \item b) \textbf{Đúng}.
    \item c) \textbf{Sai}.
    \item d) \textbf{Đúng}.
\end{itemize}
}
\end{ex}

\begin{ex} %[3D1V08]
Một cửa hàng có 36 quả bóng xếp đều vào 4 túi. Xét tính đúng sai:
\choiceTF
{\True Mỗi túi có 9 quả bóng.}
{\True Nếu bán đi 2 túi thì đã bán 18 quả bóng.}
{Số bóng còn lại sau khi bán 2 túi là 15 quả.}
{\True Số bóng còn lại sau khi bán 2 túi là 18 quả.}
\loigiai{
\begin{itemize}
    \item Mỗi túi: $36 : 4 = 9$ quả: Đúng.
    \item 2 túi bán: $9 \times 2 = 18$ quả: Đúng.
    \item Còn lại: $36 - 18 = 18$ quả $\neq 15$ quả. Ý 3 Sai, Ý 4 Đúng.
\end{itemize}
\begin{itemize}
    \item a) \textbf{Đúng}.
    \item b) \textbf{Đúng}.
    \item c) \textbf{Sai}.
    \item d) \textbf{Đúng}.
\end{itemize}
}
\end{ex}

\textbf{Phần 3. Trả lời ngắn}
\setcounterpageref{ex}{0}

\begin{ex} %[3D1H08]
Tính tổng: $456 + 132$.
\shortans{588}
\loigiai{
\textbf{Đáp số:} $588$\\
$456 + 132 = 588$.
}
\end{ex}

\begin{ex} %[3D1H08]
Tính tích: $4 \times 9$.
\shortans{36}
\loigiai{
\textbf{Đáp số:} $36$\\
$4 \times 9 = 36$.
}
\end{ex}

\begin{ex} %[3D1V08]
Tính thương: $28 : 4$.
\shortans{7}
\loigiai{
\textbf{Đáp số:} $7$\\
$28 : 4 = 7$.
}
\end{ex}

\begin{ex} %[3D1V08]
Tìm $x$ biết: $x - 150 = 350$.
\shortans{500}
\loigiai{
\textbf{Đáp số:} $500$\\
$x = 350 + 150 = 500$.
}
\end{ex}

\begin{ex} %[3D1V08]
Tính chu vi hình tam giác có các cạnh dài $15\text{ cm}, 20\text{ cm}, 25\text{ cm}$.
\shortans{60}
\loigiai{
\textbf{Đáp số:} $60$\\
$15 + 20 + 25 = 60\text{ cm}$.
}
\end{ex}

\begin{ex} %[3D1C08]
Tính nhanh: $2 \times 3 \times 4$.
\shortans{24}
\loigiai{
\textbf{Đáp số:} $24$\\
$6 \times 4 = 24$.
}
\end{ex}
